% =========================
% FILE: main.tex  (v1.0.0 – pre-Ipcha-review)
% PURPOSE: Submission-ready English IMRaD article (article class)
% COMPILE: pdflatex main.tex && bibtex main && pdflatex main.tex && pdflatex main.tex
% =========================
\documentclass[11pt,a4paper]{article}

% ---------- Language / typography ----------
\usepackage[T1]{fontenc}
\usepackage[utf8]{inputenc}
\usepackage{lmodern}
\usepackage[english]{babel}
\usepackage{csquotes}
\usepackage{microtype}

% ---------- Layout ----------
\usepackage[a4paper,margin=2.5cm]{geometry}

% ---------- Math / tables / graphics ----------
\usepackage{amsmath,amssymb,mathtools}
\usepackage{booktabs}
\usepackage{graphicx}

% ---------- Algorithms ----------
\usepackage{algorithm}
\usepackage{algpseudocode}

% ---------- References / links ----------
\usepackage{hyperref}
\usepackage{url}
\hypersetup{
  colorlinks=true,
  linkcolor=black,
  citecolor=black,
  urlcolor=black,
  pdftitle={Ipcha Mistabra -- Institutionalized Contradiction as a Verification Primitive for Epistemically Robust Text Reviews},
  pdfauthor={Oliver Baer}
}

\newcommand{\methodname}{\textsc{Ipcha Mistabra}}
\newcommand{\protshort}{\textsc{IM~Protocol}}
\newcommand{\version}{v1.0.0}
\newcommand{\doiPlaceholder}{10.5281/zenodo.\textbf{XXXXXXX}}

\title{\textbf{Ipcha Mistabra -- Institutionalized Contradiction as a Verification Primitive for Epistemically Robust Text Reviews}\\[4pt]
\large A Trialectic Questioning Protocol with Reproducible Artifacts and IP-Aware Publication}
\author{Oliver Baer\\nyxCore Systems Research Division\\\texttt{research@nyxcore.cloud}}
\date{March 2026}

\begin{document}
\maketitle

% ------------------------------------------------
\section*{Executive Summary}
% ------------------------------------------------
Safety-critical decisions frequently depend on textual artifacts---specifications, policies, architecture documents, and threat models.
Yet text-based reviews, whether manual or LLM-assisted, remain susceptible to confirmation bias, hallucination, and insufficient self-validation.
This paper introduces \methodname{}, a general, domain-agnostic verification protocol that institutionalizes contradiction as a mandatory review phase:
a \emph{Proponent} extracts claims and produces an initial assessment; an \emph{Ipcha Agent} (contradictor) systematically constructs counter-hypotheses and demands evidence; an \emph{Auditor} consolidates findings, enforces evidence binding, and renders release-gate decisions.
We define formal objects (Claim, Evidence, Finding), the \emph{Ipcha Score} as a divergence-based measure of epistemic correction, provide implementable pseudocode, and outline a reproducible experiment design (metrics, statistical analysis plan).
The manuscript integrates artifact and IP constraints for Zenodo/DOI publication and double-blind submissions (placeholders for data/results).

% ------------------------------------------------
\begin{abstract}
\textbf{Background:}
Reviews of security and compliance text artifacts are central to established frameworks (e.g., BSI IT-Grundschutz, NIST SSDF), yet remain heterogeneous and error-prone in practice~\cite{bsi_standard_100_2,nist_sp800_218}.
\\
\textbf{Problem:}
Modern text-based assistants can over-confirm user assumptions (sycophancy) and, without independent counter-review, contribute to a false sense of security~\cite{sharma_sycophancy}.
\\
\textbf{Objective:}
We introduce a general methodology that operationalizes contradiction as a mandatory verification stage, enabling systematic falsification or evidence-based support of claims and assumptions in text.
\\
\textbf{Method:}
\methodname{} defines a trialectic workflow comprising \emph{Proponent} (thesis), \emph{Ipcha Agent} (antithesis), and \emph{Auditor} (synthesis).
The protocol extracts atomic claims, generates adversarial challenge questions and counter-hypotheses, requires evidence, and produces auditable findings and remediation plans.
\\
\textbf{Evaluation (plan):}
We outline a within-subject experiment comparing a single-pass baseline with \methodname{} across multiple artifact classes (placeholder datasets), with metrics (recall/precision, risk-weighted residual errors, effort), inter-rater reliability, and a statistical analysis plan (paired tests, mixed models, FDR control).
\\
\textbf{Expected results \& limitations:}
Higher finding coverage and lower residual errors at increased effort; limitations include ground-truth cost, domain specifics, and evidence-dependent validation.
\\
\textbf{Reproducibility/IP:}
The paper contains a Zenodo-ready artifact manifest and summarizes IP constraints (DE/EU/US) as well as Zenodo metadata publicity~\cite{zenodo_policies,patg_par3,epue_art54}.
\end{abstract}

\vspace{0.5em}
\noindent\textbf{Keywords:}
security analysis; hardening; text review; adversarial questioning; reproducibility; research artifacts; IP strategy; Ipcha Score.

\vspace{0.5em}
\noindent\textbf{ACM CCS Concepts (suggested):}
\begin{itemize}
  \item Security and privacy $\rightarrow$ Software security engineering.
  \item Software and its engineering $\rightarrow$ Software verification and validation.
  \item Computing methodologies $\rightarrow$ Natural language processing; Artificial intelligence.
\end{itemize}

% ================================================================
\section{Introduction}
% ================================================================
Security-relevant defects often originate early: in requirements, architecture texts, policies, and threat models.
Frameworks such as BSI IT-Grundschutz and NIST SSDF emphasize repeatable, documented security practices including reviews and evidence trails~\cite{bsi_standard_100_2,nist_sp800_218}.
In practice, however, text reviews are inconsistent: reviewers tend to focus on plausibility rather than falsification, and LLM-assisted reviews can confirm user assumptions even when incorrect (sycophancy)~\cite{sharma_sycophancy}.

Most mitigation strategies target training-time or inference-time behavioral alignment.
These approaches improve average behavior but leave a deeper architectural problem unaddressed: the same system instance is often responsible for both proposing and implicitly validating its own reasoning.

\subsection{Research Gap}
There is no general, \emph{implementable}, and \emph{artifact-ready} protocol that
(i)~extracts claims from text artifacts,
(ii)~operationalizes contradiction as a mandatory step,
(iii)~enforces evidence binding, and
(iv)~produces an auditable finding and remediation log,
compatible with double-blind submissions and artifact policies at security venues~\cite{usenix_sec26_cfp,usenix_sec26_artifacts,ccs_2026_cfp}.

\subsection{Contributions}
\begin{enumerate}
  \item \textbf{Protocol:} \methodname{} as a trialectic verification primitive for text artifacts.
  \item \textbf{Formalization:} Definitions (claim, assumption, evidence, finding) and pseudocode.
  \item \textbf{Ipcha Score:} A divergence-based quantitative signal for epistemic correction.
  \item \textbf{Evaluation (plan):} Measurement concept and statistical analysis plan (placeholders for data/results).
  \item \textbf{Reproducibility \& IP:} Artifact manifest, Zenodo-ready metadata, and IP constraints~\cite{zenodo_policies,patg_par3,geschgehg_par2,epue_art54}.
\end{enumerate}

% ================================================================
\section{Related Work}
% ================================================================

\subsection{Document-Based Security Processes}
BSI IT-Grundschutz describes a systematic approach to information security in which documentation and security review play central roles~\cite{bsi_standard_100_2}.
NIST SP~800-218 (SSDF) specifies high-level practices for reducing vulnerability risk across the SDLC~\cite{nist_sp800_218}.
These standards motivate reviews but do not define a domain-agnostic, contradiction-institutionalizing text-review primitive.

\subsection{Failure Modes of LLM-Assisted Text Analysis}
Sycophancy denotes the tendency of models to confirm user opinions even at the expense of truth~\cite{sharma_sycophancy}.
For safety-relevant text reviews this is especially critical, because ``confirmation'' can be mistakenly interpreted as ``validation.''

\subsection{Adversarial Review and Debate Setups}
Multi-instance or debate-based approaches can improve quality and factuality by forcing explicit counter-arguments~\cite{du_multiagent_debate}.
\methodname{} operationalizes this idea as a verifiable protocol (roles, artifacts, gate rules) rather than a purely prompt-based heuristic.

\subsection{Reproducibility and Artifact Policies}
Security venues increasingly require and formally evaluate artifacts.
USENIX Security~2026 requires open artifact sharing \emph{at submission time} (with justified exceptions) and ties acceptance to continued availability~\cite{usenix_sec26_artifacts,usenix_sec26_cfp}.
ACM CCS~2026 requires that artifacts be accessible during double-blind review (anonymous URLs/credentials)~\cite{ccs_2026_cfp}.
ACM additionally defines standardized artifact badges~\cite{acm_artifact_badging}.

% ================================================================
\section{Method}
% ================================================================

\subsection{Problem Formulation}
Given a text artifact $T$ (e.g., specification, policy) and context $K$ (domain, threat model, norms/authority documents).
The objective is to produce $R$, an auditable report containing findings~$F$, evidence bindings, and a remediation plan; optionally with a gate decision (release/block).

\subsection{Formal Definitions}

\textbf{Definition~1 (Claim).}
A \emph{claim} $c$ is an atomic, verifiable statement from $T$ (e.g., ``data is encrypted in transit'').
The set of all claims is $C=\{c_1,\dots,c_n\}$.

\textbf{Definition~2 (Assumption).}
An \emph{assumption} $a$ is a prerequisite for $c$ that is not explicitly evidenced in $T$ (e.g., ``a valid certificate chain is in place'').
$A=\{a_1,\dots,a_m\}$.

\textbf{Definition~3 (Evidence).}
\emph{Evidence} $e$ is a traceable attestation (text passage, standard section, test log, code/configuration proof) supporting or refuting a claim.
$E=\{e_1,\dots,e_k\}$.

\textbf{Definition~4 (Finding).}
A finding is a tuple
\[
  f = (\mathit{id}, c, \mathit{sev}, \mathit{rationale}, E_f, \mathit{remediation}, \mathit{status})
\]
with evidence subset $E_f \subseteq E$, severity $\mathit{sev}$, and status $\in \{\text{accepted}, \text{refuted}, \text{needs-evidence}\}$.

\textbf{Definition~5 (Gate Rule).}
A gate rule $\mathcal{G}(F)\in\{0,1\}$ decides whether release/approval is permitted, e.g.:
\[
  \mathcal{G}(F)=1 \iff \nexists\, f\in F:\;(f.\mathit{status}\neq\text{refuted}) \;\land\; (f.\mathit{sev} \ge \text{high}).
\]

\subsection{Protocol Overview (Trialectic)}
\methodname{} employs three roles (see Figure~\ref{fig:architecture}):
\begin{itemize}
  \item \textbf{Proponent~(P):} extracts claims/assumptions, produces initial review $O_p$ and initial findings $F_p$.
  \item \textbf{Ipcha Agent~(I):} generates strongest counter-narratives, counter-hypotheses, and evidence demands; produces $O_i$ and $F_i$.
  \item \textbf{Auditor~(A):} consolidates, deduplicates, prioritizes; enforces evidence; renders gate decision; produces final $O_f$.
\end{itemize}

\begin{figure}[t]
  \centering
  \includegraphics[width=\linewidth]{figures/ipcha_architecture_paper.pdf}
  \caption{High-level architecture of the \methodname{} protocol: reasoning layer (Proponent $\rightarrow$ Ipcha Agent $\rightarrow$ Auditor) with a verification and memory layer providing authority grounding, epistemic memory, and workflow persistence.}
  \label{fig:architecture}
\end{figure}

\subsection{Step-by-Step Protocol}

\paragraph{Step~0: Scope Fixation.}
Define objective, artifact class, threat model, and authority documents (e.g., BSI/NIST references)~\cite{bsi_standard_100_2,nist_sp800_218}.

\paragraph{Step~1: Claim Extraction.}
Segment $T$; extract $C$ and $A$; normalize to atomic statements.

\paragraph{Step~2: Proponent Pass.}
Generate $O_p$ and findings $F_p$ with evidence requirements.

\paragraph{Step~3: Ipcha Pass (Contradiction).}
Falsify: counterexamples, contradictions, missing prerequisites, norm/policy conflicts; generate $O_i$, $F_i$.

\paragraph{Step~4: Audit Resolution.}
Merge $F=\mathrm{merge}(F_p, F_i)$; status decisions; remediation plan.

\paragraph{Step~5: Gate.}
Apply $\mathcal{G}$ (release/block).

\paragraph{Step~6: Metrics.}
Compute review metrics (e.g., finding counts, effort; Ipcha Score as divergence indicator).

\subsection{Pseudocode}
\begin{algorithm}[h]
\caption{\methodname{}: Institutionalized Contradiction for Text Artifact Review}
\begin{algorithmic}[1]
\Require Text artifact $T$, context $K$, gate rule $\mathcal{G}$
\Ensure Report $R$, findings $F$, gate decision $g$
\State $(C,A) \gets \textsc{ExtractClaims}(T)$
\State $(O_p, F_p) \gets \textsc{ProponentReview}(T,K,C,A)$
\State $(O_i, F_i) \gets \textsc{IpchaContradict}(T,K,C,A,O_p,F_p)$
\State $F \gets \textsc{MergeAndNormalize}(F_p,F_i)$
\State $(O_f, F) \gets \textsc{AuditAndResolve}(T,K,O_p,O_i,F)$
\State $g \gets \mathcal{G}(F)$
\State $R \gets \textsc{AssembleReport}(T,K,O_f,F,g)$
\State \Return $R, F, g$
\end{algorithmic}
\end{algorithm}

% ================================================================
\section{Ipcha Score}
% ================================================================

We define the \textbf{Ipcha Score} as a divergence-based measure of epistemic correction between the Proponent output and the final Auditor output:

\begin{equation}
  \mathit{IS} \;=\; 1 - \cos\!\bigl(\mathrm{embed}(A_p),\;\mathrm{embed}(A_f)\bigr)
  \label{eq:ipcha-score}
\end{equation}

\noindent where $A_p$ is the Proponent output, $A_f$ is the final Auditor output, and $\mathrm{embed}(\cdot)$ denotes a shared embedding function.

\paragraph{Interpretation (see Figure~\ref{fig:score}).}
\begin{itemize}
  \item \textbf{Low score} ($< 0.15$): the original answer was already robust, or critique was weak.
  \item \textbf{Mid-range} ($0.15$--$0.45$): normal dialectical refinement; the expected operating range for useful adversarial review.
  \item \textbf{High score} ($> 0.45$): substantial correction, indicating that important weaknesses were discovered and resolved.
\end{itemize}

\begin{figure}[t]
  \centering
  \includegraphics[width=0.82\linewidth]{figures/ipcha_score_chart.pdf}
  \caption{Conceptual interpretation of the Ipcha Score. Higher values indicate stronger correction pressure between Proponent and final validated output.}
  \label{fig:score}
\end{figure}

% ================================================================
\section{Evaluation}
% ================================================================

\subsection{Research Questions}
\begin{itemize}
  \item \textbf{RQ1:} Does \methodname{} increase finding recall/precision relative to single-pass reviews?
  \item \textbf{RQ2:} Does \methodname{} reduce risk-weighted residual errors and policy violations?
  \item \textbf{RQ3:} What cost/benefit trade-offs arise (time, compute/tokens, human rework)?
\end{itemize}

\subsection{Design (Placeholder-Driven)}
\textbf{Design:} Within-subject: each artifact is reviewed with baseline (single-pass) and \methodname{}.\\
\textbf{Artifact classes (placeholder):} $D_{\text{spec}}$, $D_{\text{policy}}$, $D_{\text{threat}}$, $D_{\text{compliance}}$.\\
\textbf{Ground truth:} Expert panel or consensus labeling; IRR is reported (see below).

\subsection{Metrics}
\begin{table}[h]
\centering
\begin{tabular}{@{}ll@{}}
\toprule
Metric & Definition / Measurement \\ \midrule
Finding Recall & Proportion of confirmed ground-truth findings discovered \\
Finding Precision & Proportion of reported findings confirmed by experts \\
Risk-Weighted Residual Error & Sum of severity for open/high findings after gate \\
Policy/Compliance Violations & Rate of violations against defined rules/standards \\
Review Effort & Time (min), compute/tokens (placeholder), rework (min) \\
Ipcha Score & Divergence measure per Equation~\ref{eq:ipcha-score} \\
IRR & Cohen's $\kappa$ / Krippendorff's $\alpha$ (depending on setup) \\ \bottomrule
\end{tabular}
\caption{Evaluation metrics (numeric values: placeholder).}
\end{table}

\subsection{Illustrative Benchmark Structure}
\begin{table}[h]
\centering
\small
\begin{tabular}{lccc}
\toprule
Metric & Baseline & \protshort{} & Improvement \\
\midrule
Error rate & 18.5\% & 2.1\% & $-88$\% \\
Compliance violations & 11.2\% & 0.9\% & $-92$\% \\
Human review effort & 100\% & 16\% & $-84$\% \\
Output validity & 72.0\% & 96.5\% & $+24.5$\,pp \\
\bottomrule
\end{tabular}
\caption{Illustrative benchmark structure for comparing baseline single-pass review and the \protshort{}. Values are placeholders for demonstration.}
\label{tab:bench}
\end{table}

\subsection{Statistical Analysis Plan}
\textbf{Paired comparisons:} Wilcoxon signed-rank (nonparametric) or paired $t$-test (if normality is plausible), plus effect size (Cliff's delta / Cohen's $d$).\\
\textbf{Mixed models:} Logistic/mixed effects for violations; ``artifact'' as random effect.\\
\textbf{Multiple testing:} FDR control (Benjamini--Hochberg).\\
\textbf{Power:} A~priori power analysis (placeholder; depends on expected effect size).

% ================================================================
\section{Discussion}
% ================================================================

\subsection{Institutionalized Contradiction as Systems Primitive}
The \protshort{} reframes AI reliability as a systems problem rather than a pure model-alignment problem.
This is analogous to fault tolerance in distributed systems: robustness comes from structured checks between components, not from assuming any one component is perfect.
A key advantage is that dissent becomes infrastructure---critique is no longer dependent on user prompt quality, developer vigilance, or informal red-teaming, but an expected and logged stage of the workflow.

\subsection{Reproducibility and Artifact Compliance}
USENIX Security~2026 requires open artifact sharing at submission; acceptance is tied to continued availability~\cite{usenix_sec26_artifacts,usenix_sec26_cfp}.
CCS~2026 requires access to artifacts during double-blind review (anonymous URLs/credentials)~\cite{ccs_2026_cfp}.
The artifact manifest in Appendix~A is designed to satisfy both.

\subsection{IP Constraints (DE/EU/Zenodo/Trade Secrets)}
Under German law, the state of the art includes everything made publicly available before the relevant date and is novelty-destroying~\cite{patg_par3}.
The EPC defines novelty accordingly (Art.~54)~\cite{epue_art54}.
If parts of the approach are to be treated as trade secrets, GeschGehG requires, among other things, secrecy and appropriate protective measures~\cite{geschgehg_par2}.
Zenodo states that metadata are CC0-licensed and publicly harvestable; metadata are therefore not a suitable repository for secrets~\cite{zenodo_policies}.
For computer-implemented aspects, note that computer programs ``as such'' are not patentable; however, computer-implemented inventions may be patentable if they solve a technical problem~\cite{dpma_cie}.

\paragraph{IP Note for USENIX~2026.}
If you wish to preserve DE/EU patent options, ``open artifacts at submission time'' may conflict with ``file before you publish.''
This tension is one reason the manuscript contains IP sections and redaction strategies~\cite{usenix_sec26_artifacts,patg_par3,epue_art54}.

% ================================================================
\section{Conclusion and Outlook}
% ================================================================
\methodname{} institutionalizes contradiction as a verification primitive for safety-relevant text artifacts.
We provide a formal specification, pseudocode, the Ipcha Score as a quantitative divergence signal, an evaluation design, and an artifact manifest for reproducible research.
We argue that institutionalized contradiction should be treated as a foundational primitive for enterprise AI workflows, especially where decisions must survive scrutiny beyond surface-level fluency.

Future work includes
(i)~standardized benchmarks,
(ii)~empirical field studies,
(iii)~robustness analysis against adversarial text inputs, and
(iv)~tooling for semi-automated evidence collection.

\section*{Acknowledgments}
% Add acknowledgments here.

\section*{Funding}
No external funding.

\section*{Author Contributions}
O.~Baer: Conceptualization; Methodology; Software; Validation; Writing (Original Draft); Writing (Review \& Editing).

\section*{Conflict of Interest}
The author declares no conflict of interest.

\bibliographystyle{unsrt}
\bibliography{references}

% ================================================================
\appendix
% ================================================================

\section{Artifact Package (Zenodo-Ready) and File Manifest}

\begin{table}[h]
\centering
\begin{tabular}{@{}llp{7.0cm}@{}}
\toprule
File & Format & Purpose \\ \midrule
paper\_ipcha\_mistabra\_\version.pdf & PDF & Main paper (citable), contains protocol \\
ipcha\_mistabra\_zenodo\_bundle\_\version.zip & ZIP & Release bundle (folder structure, artifacts, docs) \\
SHA256SUMS\_\version.txt & TXT & Checksums (integrity; ZIP + PDF) \\
README.md & MD & Quickstart, citation note (DOI) \\
REPRODUCIBILITY.md & MD & Reproduction instructions, environment, expected outputs \\
MANIFEST.json & JSON & Machine-readable table of contents/metadata \\
LICENSE\_DOCS & TXT & Document license (e.g., CC~BY~4.0 note) \\
LICENSE\_CODE & TXT & Code license (MIT/Apache-2.0 etc.) \\
CITATION.cff & YAML & Citation hints (DOI placeholder) \\
CHANGELOG.md & MD & Versioning (SemVer), release notes \\
SECURITY.md & MD & Vulnerability/issue reporting \\
PRIOR\_ART\_LOG.md & MD & State-of-the-art / related-work log \\
NDA\_template.md & MD & (Optional) NDA template for restricted collaboration \\
invention\_dossier\_confidential.pdf & PDF & (Typically not public) IP dossier \\
invention\_dossier\_public\_redacted.pdf & PDF & (Optional) redacted, public version \\ \bottomrule
\end{tabular}
\caption{Artifact manifest (suggested file names; contents: placeholder).}
\end{table}

\section{Reproducibility: Commands, Dockerfile Skeleton, Packaging}

\subsection*{Dockerfile Skeleton (Placeholder)}
\begin{verbatim}
# artifacts/tooling/Dockerfile
FROM python:3.12-slim
WORKDIR /app
COPY requirements.txt .
RUN pip install --no-cache-dir -r requirements.txt
COPY src/ ./src/
COPY run.sh .
ENTRYPOINT ["bash","./run.sh"]
\end{verbatim}

\subsection*{Example CLI}
\begin{verbatim}
unzip ipcha_mistabra_zenodo_bundle_v1.0.0.zip
cd ipcha_mistabra_release_v1.0.0/artifacts/tooling
docker build -t im-tooling:v1 .
docker run --rm -v "$(pwd)/../examples:/data" im-tooling:v1 \
  --input /data/example_inputs/spec.txt \
  --context /data/example_inputs/context.json \
  --output /data/example_outputs/report.json
\end{verbatim}

\subsection*{Packaging and SHA256}
\begin{verbatim}
zip -r ipcha_mistabra_zenodo_bundle_v1.0.0.zip \
  ipcha_mistabra_release_v1.0.0/
sha256sum paper_ipcha_mistabra_v1.0.0.pdf \
  ipcha_mistabra_zenodo_bundle_v1.0.0.zip \
  > SHA256SUMS_v1.0.0.txt
\end{verbatim}

\section{Zenodo Metadata Snippets (Copy-Paste)}

\subsection*{Title}
\begin{verbatim}
Ipcha Mistabra: A General Method for Structured Text Review
and Security-Oriented Hardening (v1.0.0)
\end{verbatim}

\subsection*{Description (Short Abstract)}
\begin{verbatim}
Ipcha Mistabra is a general, domain-agnostic method for
structured text review (questioning) that extracts claims and
assumptions from text artifacts, adversarially challenges them,
and produces evidence-based findings and remediations.
The deposit contains the paper, artifacts/examples, checksums,
and reproducibility instructions.
\end{verbatim}

\subsection*{Keywords}
\begin{verbatim}
security analysis; hardening; text review; adversarial
questioning; reproducibility; artifacts; IP; Ipcha Score
\end{verbatim}

\subsection*{License}
\begin{verbatim}
Documents: CC BY 4.0
Code/Tooling: MIT or Apache-2.0
\end{verbatim}

\end{document}
